% ============================================================
% CHƯƠNG 5: KẾT LUẬN VÀ HƯỚNG PHÁT TRIỂN
% ============================================================
\clearpage
\section{Kết luận và hướng phát triển}

% ============================================================
\subsection{Kết luận}

\subsubsection{Đánh giá mức hoàn thành so với mục tiêu}

Đề tài đặt ra ba mục tiêu ban đầu; bảng dưới đây đối chiếu
từng mục tiêu với kết quả đạt được:

\begin{enumerate}
    \item \textbf{Data Warehouse chuẩn hóa.}
    Kho dữ liệu SQLite theo Star Schema đã được xây dựng hoàn chỉnh:
    1 bảng sự kiện (\texttt{fact\_flights}, $\approx$5{,}3M bản ghi),
    4 bảng chiều và 6 chỉ mục OLAP.
    ETL pipeline streaming kết hợp tối ưu dtype (\texttt{float32}/\texttt{int8})
    giữ RAM ổn định $\approx$120\,MB/chunk, hoàn thành trong 3--8 phút ---
    chứng minh kiến trúc phân lô hoạt động hiệu quả trên máy cá nhân.

    \item \textbf{Dashboard OLAP đa chiều.}
    Dashboard Streamlit 6 trang (Tổng quan, Tuyến đường, Thời gian,
    Nguyên nhân, Dự báo, Explorer) kết hợp Plotly,
    hiện thực đầy đủ các phép OLAP (Slice, Dice, Roll-up, Drill-down)
    thông qua bộ lọc tương tác theo tháng/quý, hãng và địa lý.
    Mỗi trang tập trung một góc nhìn phân tích riêng,
    giúp người dùng khám phá dữ liệu một cách có hệ thống.

    \item \textbf{Mô hình dự báo delay.}
    Framework so sánh 5 thuật toán (Logistic Regression, Random Forest,
    Neural Network, XGBoost, LightGBM) trên cùng tập dữ liệu và 7 chỉ số.
    Mô hình tốt nhất (XGBoost) đạt AUC-ROC $\approx$ 0{,}75 ---
    vượt đáng kể so với baseline ngẫu nhiên (0{,}5) --- với bộ feature
    hoàn toàn ``dự báo''.
\end{enumerate}

% ------------------------------------------------------------
\subsubsection{Insight thực tế và khuyến nghị từ dữ liệu 2015}

Kết hợp phân tích OLAP trên dashboard và Feature Importance từ mô hình,
nhóm rút ra các quy luật quan trọng, kèm \textbf{khuyến nghị hành động}
cho các bên liên quan:

\begin{itemize}
    \item \textbf{Hiệu ứng dây chuyền là nguyên nhân delay lớn nhất.}
          Late Aircraft và NAS chiếm tỷ trọng trễ cao nhất ---
          hệ thống hàng không rất nhạy với gián đoạn ban đầu:
          một chuyến delay buổi sáng có thể lan sang 3--5 chuyến tiếp theo.
          \textit{Khuyến nghị cho hãng:} tăng buffer thời gian quay đầu tàu bay
          (\textit{turnaround time}) cho các chuyến nối tiếp,
          đặc biệt tại sân bay hub.

    \item \textbf{Mùa hè (tháng 6--7) và mùa đông (tháng 12--1) là hai đỉnh delay.}
          Trùng với mùa bão và mùa tuyết --- hai giai đoạn mà tỷ lệ trễ
          vượt trung bình năm đáng kể.
          \textit{Khuyến nghị cho quản lý sân bay:} chuẩn bị phương án
          dự phòng vận hành riêng cho hai mùa cao điểm,
          bao gồm bổ sung nhân lực mặt đất và tối ưu lịch bay.

    \item \textbf{Chuyến sáng sớm (trước 8h) ít delay nhất.}
          Chưa chịu ảnh hưởng tích lũy từ các trễ trong ngày.
          \textit{Khuyến nghị cho hành khách:} ưu tiên chuyến sáng sớm
          nếu cần đảm bảo đúng giờ cho kết nối quan trọng.

    \item \textbf{Hãng bay ảnh hưởng có hệ thống đến đúng giờ.}
          Chênh lệch OTP giữa hãng tốt nhất và kém nhất lên đến 15--20 điểm phần trăm,
          phản ánh khác biệt trong chiến lược quay đầu, bảo trì và xếp lịch.
          \textit{Khuyến nghị cho ngành:} benchmark OTP công khai
          tạo áp lực cải thiện từ phía hãng có tỷ lệ trễ cao.

    \item \textbf{Quãng đường và thời gian bay là feature quan trọng nhất trong mô hình.}
          Chuyến bay dài chịu nhiều rủi ro hơn (thời tiết, không phận, tiêu thụ nhiên liệu).
          Yếu tố địa lý (bang sân bay) cũng được mô hình khai thác mạnh,
          cho thấy hạ tầng và khí hậu vùng miền là nhân tố phân biệt rõ.
\end{itemize}

% ============================================================
\subsection{Hạn chế}

\begin{enumerate}
    \item \textbf{Dữ liệu chỉ bao phủ năm 2015.}
    Insight và mô hình phản ánh thực trạng một năm cụ thể.
    Sự thay đổi về cơ sở hạ tầng, quy định FAA,
    và hành vi di chuyển sau đại dịch COVID-19
    có thể làm mô hình không còn phù hợp với thực tế hiện tại.
    Tuy nhiên, \textit{phương pháp luận} (kiến trúc DW, pipeline so sánh mô hình)
    vẫn áp dụng được khi thay đổi nguồn dữ liệu.

    \item \textbf{SQLite không phù hợp cho production.}
    SQLite là single-writer, không hỗ trợ concurrent access
    và không có cơ chế phân tán.
    Một số OLAP query phức tạp mất vài giây ---
    chấp nhận được cho đồ án nhưng không đủ cho hệ thống thực
    phục vụ nhiều người dùng đồng thời.

    \item \textbf{AUC-ROC $\approx$ 0{,}75 --- còn khoảng cách so với mô hình lý tưởng.}
    Mô hình chủ động \textbf{không dùng \texttt{departure\_delay}}
    để đảm bảo tính dự báo (trước khi cất cánh).
    Đổi lại, AUC bị giới hạn vì thiếu dữ liệu thời tiết thực tế,
    thông tin bảo trì tàu bay, và lịch sử delay từng tail number ---
    các yếu tố có thể cải thiện AUC thêm 0{,}03--0{,}10.

    \item \textbf{Chưa có Clustering / Association Rules.}
    Đề tài tập trung vào classification (dự báo delay)
    và OLAP (phân tích đa chiều).
    Các kỹ thuật khai phá khác như phân cụm sân bay theo hành vi delay
    hoặc luật kết hợp (``hãng A + sân bay ORD + tháng 12 $\Rightarrow$ delay'')
    chưa được thực hiện --- có thể là hướng bổ sung.
\end{enumerate}

% ============================================================
\subsection{Hướng phát triển}

\begin{enumerate}
    \item \textbf{Mở rộng dữ liệu nhiều năm.}
    Tích hợp dữ liệu 2010--2023 từ Bureau of Transportation Statistics (BTS)
    để phát hiện xu hướng dài hạn, đánh giá ảnh hưởng COVID-19 (2020--2021),
    và huấn luyện mô hình tổng quát hơn trên phân phối đa dạng.

    \item \textbf{Tích hợp dữ liệu thời tiết.}
    Kết hợp dữ liệu NOAA (nhiệt độ, lượng mưa, tốc độ gió)
    theo sân bay và ngày --- đây là yếu tố delay lớn nhất
    mà mô hình hiện tại chưa có.

    \item \textbf{Nâng cấp hạ tầng lưu trữ.}
    \begin{itemize}
        \item Thay SQLite bằng \textbf{PostgreSQL} hoặc \textbf{ClickHouse}
              để hỗ trợ concurrent access và OLAP query nhanh hơn.
        \item Containerize hệ thống bằng \textbf{Docker},
              deploy dashboard lên cloud (AWS EC2 / Google Cloud Run).
        \item Xây dựng \textbf{REST API} (FastAPI) cho dự báo delay
              theo yêu cầu thời gian thực.
    \end{itemize}

    \item \textbf{Giải thích mô hình với SHAP.}
    Áp dụng SHAP (SHapley Additive exPlanations) \cite{lundberg2017shap}
    để giải thích \textit{từng dự báo cụ thể}:
    tại sao chuyến bay này bị dự báo delay? Feature nào đóng góp nhiều nhất?
    Dashboard có thể hiển thị không chỉ ``có delay'' mà còn ``tại sao delay''.

    \item \textbf{Pipeline streaming thời gian thực.}
    Sử dụng Apache Kafka nhận dữ liệu chuyến bay real-time từ API hàng không,
    kết hợp Spark Streaming để dự báo delay ngay khi chuyến bay được đặt lịch ---
    chuyển hệ thống từ phân tích lịch sử sang hỗ trợ quyết định tức thời.
\end{enumerate}
